La République du Bénin a engagé la mise en place d'un système national
de vidéoprotection dont la vidéo-verbalisation par vision par ordinateur
est une application clef. Le premier maillon du pipeline, la détection
des véhicules, échoue le plus souvent sur les véhicules éloignés, de
petite taille à l'image ou masqués dans les embouteillages, qui sont
précisément les configurations dominantes du trafic urbain
ouest-africain. Nous
étudions si un fine-tuning de YOLOv8n sur du trafic urbain dense
améliore cette détection. Faute de jeu de données béninois public, nous
utilisons un sous-ensemble reproductible de 3\,000 images
(25\,688 véhicules annotés) de BMD-45, jeu de caméras CCTV de Bengaluru
structurellement comparable (caméras fixes, forte densité, $56\,\%$ de
deux-roues), remappé vers quatre classes alignées sur COCO. Le protocole
de comparaison garantit l'équité entre modèle pré-entraîné et modèle
fine-tuné : mêmes images et boîtes de test, réindexation des classes,
rappel ventilé par taille d'objet selon la convention COCO avec aires
normalisées à la résolution d'entrée du réseau. Après 50 epochs, la
mAP@0,5 passe de $0{,}438$ à $0{,}825$ ($+88{,}4\,\%$) et le rappel des
objets de moins de $32 \times 32$ pixels de $0{,}122$ à $0{,}694$
($+469{,}4\,\%$), le gain étant d'autant plus fort que les objets sont
petits. L'adaptation de domaine par fine-tuning lève donc précisément le
verrou visé et quantifie le bénéfice qu'apporterait un jeu de données
local.
